% \iffalse
%<*drv>
\documentclass{ltxdoc}
\usepackage{doc}
\usepackage[T1]{fontenc}
\usepackage[latin1]{inputenc}
\parindent0pt
\begin{document}
%\OnlyDescription  %Comment out for Desription only
\setlength\hfuzz{15pt}
\sloppy
\DocInput{serial.dtx}
\end{document}
%</drv>
%\fi
% \MakeShortVerb{\|}
% \CheckSum{589}
% \section{Guide}
% \subsection{Introduction}
% This class has been designed as a local wrapper for the KOMA-Script letter 
% class to allow easy generation of serial letters.
%
% Three different approaches
% are now implemented, one based on the method described in the 
% scrlettr-guide, and one based on the textmerg package, and one using 
% the delimtxt-package. 
%
% An interface for a local configuration file is integrated, allowing for
% easy modification of the letter layout for your site.
%
% It is recommended to use the |inputenc| package with suitable options
% if interfacing to an tex-unaware database to avoid problems with characters
% not in the default character class (e.g. german Umlaut and the like).
%
% \bigskip
% TODO:
% Integration of the standard letter class. (KOMA uses some different names
% for part of the letter information).
%
% The documentation is not fully done yet.
% The examples still need
% to be translated from german to english.
%
% \bigskip
% 
% Copyright  1999 Bjoern Pedersen
%
% This program can be redistributed and/or modified under the terms
% of the LaTeX Project Public License Distributed from CTAN
% archives in directory macros/latex/base/lppl.txt; either
% version 1 of the License, or any later version.
%
%  Permission is granted to to customize the declarations in this file
%  to serve the needs of your installation.  However, no permission is
%  granted to distribute a modified version of this file under its
%  original name "serial.dtx".
%
%
%\subsection{Class options}
% \begin{enumerate}
% \item \lbrack |kein|\rbrack |vordruck| switches from preprinted to white paper 
% \lbrack default:white paper\rbrack
% \item |adrentry| use KOMA-Script based data file
% \item |merg| use textmerg style data file  \lbrack default\rbrack
% \item |delim| use delimtxt style data file  
% \item \lbrack |no|\rbrack |nofoldmarks| switches foldmarks on/off \lbrack on\rbrack
% \item |latin1/asciinew/ansi/applemac| If present, these are passed to the |inputenc|
%       package
% \item |draft| This options outputs one letter, skipping the reading
% of the data file and substituting the macros with their
% name. (useful for testing the letter)
%\end{enumerate}
%
%\subsection{Macros}
% \subsubsection{Setup and Utility Macros}
% \DescribeMacro{\MyFirstHead}|\MyFirstHead|:
%
% This generates our local header on the first page of each letter.
%
% \DescribeMacro{\MyFirstFoot}|\MyFirstFoot|:
% This generates our local header on the first page of each letter.
% Use |DeclareRobustCommand{\MyFirst[Head/Foot]}| to change it e.g. in the cfg-file.
%
% \DescribeMacro{\MyNextHead}|\MyNextHead|,
% \DescribeMacro{\MyNextFoot}|\MyNextFoot|:
%
% These are the headers/footers on the following pages and by default empty.
%
% Behind the scenes, the fancyhdr-package is used to customize these headers/
% footers.
%
% \subsubsection{KOMA-script -letter macros}
% These scrlettr-macros have to be declared either in the config-file or in the preamble.
%
% \DescribeMacro{\signature}|\signature|
%
% \DescribeMacro{\place}|\place|
%
% \DescribeMacro{\name}|\name|
%
% \DescribeMacro{\location}|\location|
%
% \DescribeMacro{\telephone}|\telephone|
%
% \DescribeMacro{\backadress}|\backadress|
%
% \subsubsection{The serial letter macros}
% \DescribeMacro{\Serial}
% The |\Serial| macro is the main interface to this class. Its semantics are
% as follows:
%
% |\Serial{adressfile}{letter-body}|
% 
% It will read an adressfile( formatted as decribed in section \ref{adfile})
% and generate the serial letters. In each version an extra counter will enumerate the letters,
% so later selected letters could be printed.
%
% The following macros will be defined in the letterbody (using the textmerg/delimtxt-options
% , these can be further customized, see below): 
%
% Entries marked with (\textbf{+}) are used internally and should always be present with this
% name if customizing the data base format. Otherwise you have to rewrite the relevant parts
% in your local class. 
%
%\DescribeMacro{\Name}|\Name|  The surname(\textbf{+}).
%
%\DescribeMacro{\VorName}|\VorName|  The first name(s)(\textbf{+}).
%
%\DescribeMacro{\Street}|\Street|  The street(\textbf{+}).
%
%\DescribeMacro{\City}|\City| The City (with ZIP-code)(\textbf{+}).
%
%\DescribeMacro{\Phone}|\Phone| The phone number.
%
%\DescribeMacro{\Geschlecht}|\Geschlecht|  The sex.
%
%\DescribeMacro{\FreeII}|\FreeII|  A multipurpose spare field.
%
%\DescribeMacro{\Comment}|\Comment| A comment field.
%
% \DescribeMacro{\Anrede}|\Anrede| Use this for the (by default german) opening. It depends on |\Geschlecht|
% being either |w| or |m| and selects \DescribeMacro{\Wanredetext}|\Wanredetext| or 
% \DescribeMacro{\Manredetext}|\Manredetext|. For other languages, customize it in the configuration file.
%
%
% \subsubsection{The Fields macro}
%\DescribeMacro{\Fields}
%  In the textmerg or delimtxt-based version, the input format can be easily customized using 
% |\renewcommand{\Fields}|. 
%
% The definition should contain the desired macro names, optionally preceeded by |+| to
% allow skipping blank lines in the input. For a full description see the |textmerg|/|delimtxt| description.  
% Remember not to omit required fields.
%
% \subsection{The Adressfile format}
% \label{adfile}
% \subsubsection{The adrentry version}
% This file with extension |.adr| should consist of lines containing:
%
%{\noindent\small |\adrentry{name}{firstname}{street}{city}{phone}{sec}{free}{comment}|}
% 
% \subsubsection{The textmerg version}
% This file with extension |.adt| has a more loosely defined format:
% Each line corresponds to one entry in the |\Fields|-macro. These blocks are repeated for each letter.
% If enabled with the |+| switch, blank lines are ignored in the input.
%
% So, for the default 
%
% |\Fields{\Name\VorName\Street\City\Phone\Geschlecht\FreeII\Comment}|
%
% there should be 8 lines for each letter in the adress file.
% \subsubsection{The delimtxt version}
% This file with extension |.del| is analogous to the textmerg file, but fields are separated by \verb+|+ .
% This delimiter can be changed by the delimtxt function |\SetDel|. Each new line starts a new letter.
% For more info on delimtxt see the documentation supplied there.
%
% \StopEventually
%
% \section{Implementation}
% \subsection{Serial-specific setup}
%<*class>
% 
% 
% Announce the classname.
%
%
%% Copyright  1999 Bjoern Pedersen
%%
%% This program can be redistributed and/or modified under the terms
%% of the LaTeX Project Public License Distributed from CTAN
%% archives in directory macros/latex/base/lppl.txt; either
%% version 1 of the License, or any later version.
%%
%%  Permission is granted to to customize the declarations in this file
%%  to serve the needs of your installation.  However, no permission is
%%  granted to distribute a modified version of this file under its
%%  original name "serial.dtx".
%    \begin{macrocode}
\def\filedate{1999/04/30}
\ProvidesClass{serial}[\filedate]
\NeedsTeXFormat{LaTeX2e}[1998/06/01]
%    \end{macrocode}
%
% Declare the necesary class options.
%
% \DescribeMacro{\firsthead}
% \DescribeMacro{\firstfoot}
% This generates our local header and footer on the first page of each letter.
% The real definition should be in the configuration file. 
%
%    \begin{macrocode}
\DeclareOption{keinvordruck}{
\AtBeginDocument{\firsthead{\MyFirstHead}
              \firstfoot{\MyFirstFoot}
%    \end{macrocode}
%
% \DescribeMacro{\nexthead} \DescribeMacro{\nextfoot}
% On the following pages, no header or footer is planned.
%
%    \begin{macrocode}

\nexthead{\MyNextHead}  
\nextfoot{\MyNextFoot}

\pagestyle{headings}
}% end of AtBeginDocument
}
%    \end{macrocode}
%
% If a preprinted paper is used, use option vordruck.
%
%    \begin{macrocode}
\DeclareOption{vordruck}%
{\AtBeginDocument{%
\firsthead{}%
\firstfoot{}%
\nexthead{}%
\nextfoot{}%
\pagestyle{empty}
}}
%    \end{macrocode}
%
% Load the required package, if any.
% To select the right behaviour, we let |\Serial| to either |\Serialx| or |\Seriala|.
% For the koma-script version make |\Fields| a no-op.
%
%    \begin{macrocode}
\newif\ifTextM
\TextMfalse
\newif\ifDelT
\DelTfalse
\newif\ifDraft
\Draftfalse

\DeclareOption{draft}{\Drafttrue}
\DeclareOption{final}{\Draftfalse}

\DeclareOption{merg}{%
\TextMtrue
        \AtEndOfClass{\global\let\MyMerge=\Merge
        \gdef\ext{.adt}
        \global\let\Serial=\Serialx}}

\DeclareOption{delim}{ 
\DelTtrue
        \AtEndOfClass{\global\let\MyMerge=\DelimRead
        \gdef\ext{.del}
        \global\let\Serial=\Serialx}}

\DeclareOption{adrentry}{
\def\Fields#1{\relax}
\countdef\LetterCounter=1
\LetterCounter=0
        \AtEndOfClass{\global\let\Serial=\Seriala}}

%    \end{macrocode}
%
% Some options for foldmarks.
%
%    \begin{macrocode}
\DeclareOption{nofoldmarks}{\AtEndOfClass{\foldmarksoff}}
\DeclareOption{foldmarks}{\AtEndOfClass{\foldmarkson}}
%    \end{macrocode}
%
% Some inputenc.sty options
%    \begin{macrocode}
\DeclareOption{ansinew}{\PassOptionsToPackage{\CurrentOption}{inputenc}}
\DeclareOption{latin1}{\PassOptionsToPackage{\CurrentOption}{inputenc}}
\DeclareOption{ascii}{\PassOptionsToPackage{\CurrentOption}{inputenc}}
\DeclareOption{applemac}{\PassOptionsToPackage{\CurrentOption}{inputenc}}
%    \end{macrocode}
%
% Setup the defaults and load |scrlettr.cls|.
%
%    \begin{macrocode}
\DeclareOption*{\PassOptionsToClass{\CurrentOption}{scrlettr}}
\ExecuteOptions{keinvordruck,foldmarks,merge}
\ProcessOptions\relax
\LoadClass[wlocfield,a4paper]{scrlettr}
%    \end{macrocode}
%
% Load required packages based on options
%
%    \begin{macrocode}
\ifTextM
\RequirePackage{textmerg}
\fi
\ifDelT
\RequirePackage{delimtxt}
\fi
%    \end{macrocode}
%
% We use a possibly very high header, so recalutate scrlettr layout ( a4paper only).
% To get those foldmarks right, we calculate the position from the top of the paper.
% |\foldhskip| adjusted for HP LJ5L. This can be done in the .cfg file.
%
%    \begin{macrocode}
\setlength{\voffset}{-5.5truemm}
\setlength{\topmargin}{-1in}
\addtolength{\topmargin}{10truemm}
\setlength{\headheight}{27truemm}
\setlength{\headsep}{4truemm}
%    \end{macrocode}    
%
% calculate vertical skip for foldmarks from other length at the beginnig
% of the document, as head height may be customized in the config file.
% This has do be done after setting backadress in the .cfg-file.
%    \begin{macrocode}
\AtEndOfClass{
\setlength{\foldvskipi}{101.5truemm}

\setlength{\foldvskipii}{47.5truemm}    % mk 960531 (war:40.5truemm)
\setlength{\foldvskipiii}{58.5truemm}    % mk 960531 (war:58.5mm)
\setlength{\foldhskip}{6truemm} % adjust for printer
%    \end{macrocode}    
%  
%    Calculate textheight from footheight.
%
%    \begin{macrocode}  
\newlength{\footheight}
\setlength{\footheight}{25truemm}
\setlength{\textheight}{297truemm}

} % End of AtEndofClass
\AtBeginDocument{
}

%    \end{macrocode}
%
% Here are the defaults for some scrlettr macros.
% They must be overriden in serial.cfg.
%

% \DescribeMacro{\signature}
% \DescribeMacro{\place}
% \DescribeMacro{\name}
% \DescribeMacro{\location}
% \DescribeMacro{\telephone}
% \DescribeMacro{\backadress}

%    \begin{macrocode}

%%------------
%%
%% this is configurable in serial.cfg
%%

 \signature{Please set this signature in your .cfg file}

 \place{Please set this place in your .cfg file}
 \name{Please set this name in your .cfg file}

 \address{Please set this address in your .cfg file }
 \location{Please set this location in your .cfg file}

 \telephone{Please set this telephone in your .cfg file}

 \backaddress{ Please set this backaddress in your .cfg file }

% \specialmail{\ }
%    \end{macrocode}
%    \begin{macrocode}
\DeclareRobustCommand{\MyFirstHead}{%
\parbox[b]{0.99\textwidth}
{  \begin{minipage}[b]{98truemm}
    {\sffamily\bfseries\fontsize{4.1truemm}{6.0truemm}\selectfont
     \ignorespaces
         SOME NAME\\
         Configure in your .cfg}%
     \end{minipage}
     \begin{minipage}[b]{48truemm}\hfill\parbox[b]{47.5truemm}
                         {here would be a graphic}
     \end{minipage} 
 }%
 }
%    \end{macrocode}
%
%    \begin{macrocode}
\DeclareRobustCommand{\MyFirstFoot}{%
\parbox[t]{\textwidth}
                {\rule{\textwidth}{.4pt}\\[0.6ex]%
                \sffamily\tiny%
You have to override this in your .cfg-file
}% end of parbox
}%
\DeclareRobustCommand{\MyNextHead}{}
\DeclareRobustCommand{\MyNextFoot}{}
%    \end{macrocode}
%
% \section{The serial letter implementation}
% \subsection{Common macros}
%
% \DescribeMacro{\Anrede}
% \DescribeMacro{\Wanredetext}
% \DescribeMacro{\Manredetext}
% In serial letters, we want to differentiate by sex , these macros will do it for us. No babel support, as for now.
% These macros can be redeclared in the configuration file.
%
%    \begin{macrocode}
\DeclareRobustCommand{\Wanredetext}{Sehr geehrte Frau}
\DeclareRobustCommand{\Manredetext}{Sehr geehrter Herr}
\DeclareRobustCommand{\Anrede}{%
\if \Geschlecht w \opening{\Wanredetext\ \Name,} \fi%
\if \Geschlecht m \opening{\Manredetext\ \Name,} \fi%
}
%    \end{macrocode}
% \DescribeMacro{\Serial}
%
% The macro |\Serial{adrfilebasename}{letterbody}| will expand to either |\Seriala| or|\Serialx| as defined 
% by the class-options.
%
% \subsection{The KOMA-Script based version}
%
% \DescribeMacro{\Seriala}
%  Here starts the |\Seriala| macro. This is a simple wrapper of the approach used in scrlettr.
% The the |komascr|-documentation for details. The support for draft modus has been added.
%
% In the example file, long lines have been split to allow inclusion in the documentation,
%  but this not required for normal use.
%    \begin{macrocode}
\DeclareRobustCommand{\Seriala}[2]{
\def\@AddFile{#1.adr}
\long\def\@body{#2}
% This is the draft part
\ifDraft
\begingroup
\global\advance\LetterCounter by1
\expandafter\def\csname Name\endcsname{\texttt{Name}}
\expandafter\def\csname VorName\endcsname{\texttt{VorName}}
\expandafter\def\csname Street\endcsname{\texttt{Street}}
\expandafter\def\csname City\endcsname{\texttt{City}}
\expandafter\def\csname Phone\endcsname{\texttt{Phone}}
\expandafter\def\csname Geschlecht\endcsname{w} % to check \Anrede
\expandafter\def\csname FreeII\endcsname{\texttt{FreeII}}
\expandafter\def\csname Comment\endcsname{\texttt{comment}}
\typeout{Addressfile: \@AddFile would be loaded.}
\begin{letter}{\VorName\ \Name\\\Street\\[1.5ex]\City}
\expandafter\noexpand\@body
\end{letter}
\endgroup%
% end of draft part
% This if the final part
\else
\begingroup
\def\adrentry##1##2##3##4##5##6##7##8{%
 \global\advance\LetterCounter by1
      \expandafter\def\csname Name\endcsname{##1}
      \expandafter\def\csname VorName\endcsname{##2}
      \expandafter\def\csname Street\endcsname{##3}
      \expandafter\def\csname City\endcsname{##4}
      \expandafter\def\csname Phone\endcsname{##5}
      \expandafter\def\csname Geschlecht\endcsname{##6}
      \expandafter\def\csname FreeII\endcsname{##7}
        \ifx\empty\FreeII
                \def\FreeII{\ }
        \fi
      \expandafter\def\csname Comment\endcsname{##8}
    \expandafter\begin{letter}{##2\ ##1\\##3\\[1.5ex]##4}
    \expandafter\noexpand\@body
    \expandafter\end{letter}
}
\InputIfFileExists{\@AddFile}
                   {\typeout{Load addressfile: \@AddFile.}}
                    {\PackageError{serial}
                                  {File \@AddFile not found}
                                  {The addressfile you wanted is not
                                   available}}
  \endgroup%
% end of final part
\fi
}
%    \end{macrocode}
%%
%%---------------
%%
%
% \subsection{The textmerg-based version}
%
% \DescribeMacro{\Serialx}
% This is the wrapper for the |textmerg|-package:
% use a textmerg style inputfile with extension |.adt|,
% containing the desired fields.
% Default is:
%
%|\Name\VorName\Street\City\Phone\Geschlecht\FreeII\Comment|
% 
% \DescribeMacro{\Fields}
% The default can be overwritten using |\Fields{fields}| before the letter. 
% Further information can be found in the documentation for the 
% textmerge-package.
% 
%    \begin{macrocode}
\Fields{\Name\VorName\Street\City\Phone\Geschlecht\FreeII\Comment}
%    \end{macrocode}
%
% The draft merge helper
% In draft mode, we parse the |\Fields| macro and define each csname
% to be its string expansion.( This should really go into |textmerg|/|delimtxt|.)
%    \begin{macrocode}
\DeclareRobustCommand{\DraftMerge}[2]{%
\typeout{The adressfile: #1 would be loaded}
\begingroup
\def\DParseFields##1{%
  \ifx##1\EndParseFields%
      \let\DNextParse\relax%
  \else%
      \let\DNextParse\DParseFields%
      \ifx##1+\relax%
      \else%
        \ifx##1-\relax%
        \else%
           \def##1{\string##1}
        \fi
      \fi
  \fi\DNextParse}%
\let\EndParseFields\DParseFields
\expandafter\DParseFields\the\Fields\EndParseFields 
\global\advance\Iteratecounter by1
\def\Geschlecht{w} % set for \Anrede -> \Wanredetext
  #2
\endgroup
}
%
\DeclareRobustCommand{\Serialx}[2]
{
\def\@AddFile{#1\ext}
\ifDraft
%begin of draft part
\DraftMerge{\@AddFile}{%
    \begin{letter}{\VorName\ \Name \\\Street\\[1.5ex]\City}%
    #2%
    \end{letter}%
}%
%end of draft part
\else
% begin of final part
\MyMerge{\@AddFile}{%
    \begin{letter}{\VorName\ \Name \\\Street\\[1.5ex]\City}%
    #2%
    \end{letter}%
}%
% end of final part
\fi
}
%    \end{macrocode}
% 
% At the end of the class we input the configuration file. This has to
% be done before we set the header and footers to allow them to be customized
% via the config file without extra macro invocations.
%
%    \begin{macrocode}
\AtEndOfClass{\InputIfFileExists{serial.cfg}%
  {\typeout{Loaded local configuration from serial.cfg}}%
  {\typeout{no local configuration in serial.cfg}
   \PackageError{serial}{Missing configuration file}
  {You should really create the %
   cfg-file as described in the documentation}}}
%
% Now recalculate the layout:
\AtEndOfClass{
\addtolength{\foldvskipi}{-1in}
\addtolength{\foldvskipi}{-\topmargin}
\addtolength{\foldvskipi}{-\headsep}
\addtolength{\foldvskipi}{-\headheight}

\addtolength{\textheight}{-1in}
\addtolength{\textheight}{-\topmargin}
\addtolength{\textheight}{-\headsep}
\addtolength{\textheight}{-\headheight}
\addtolength{\textheight}{-\footheight}
\newlength\@backheight
\newlength\@backdepth
\settoheight\@backheight{\frombackaddress}
\settodepth\@backdepth{\frombackaddress}
\addtolength{\foldvskipi}{-\@backheight}
\addtolength{\foldvskipi}{-\@backdepth}
\addtolength{\textheight}{-\@backheight}
\addtolength{\textheight}{-\@backdepth}
}
%% End of file `serial.cls'.
%    \end{macrocode}
%</class>
%
% \section{The example files}
% \subsubsection{Example letters}
%
%    \begin{macrocode} 
%<*ex1>
\documentclass[12pt,draft,adrentry]{serial}
%</ex1>
%<*ex2>
\documentclass[12pt,draft,merg]{serial}
%</ex2>
%<*ex3>
\documentclass[12pt,draft,delim]{serial}
%</ex3>
%<*exa>

\DeclareRobustCommand{\VerbX}[1]{\item[\texttt{\textbackslash#1}]}

\begin{document}

\date{\today}

\Serial{examp}{ 
\subject{Test}

\Anrede

Hier sollte der Text des Briefes stehen.
Folgende Makros stehen zur Verf�gung:
\begin{labeling}[:]{\texttt{\textbackslash Geschlecht}}
\VerbX{Name}  Ist der Nachname in der Adressendatei:\fbox{\Name}

\VerbX{VorName}  Ist der Vorname in der Adressendatei: \fbox{\VorName}

\VerbX{Street}  Ist die Strasse in der Adressendatei: \fbox{\Street} 

\VerbX{City}  Ist die Stadt in der Adressendatei: \fbox{\City} 

\VerbX{Phone} Ist die Telephonnummer in der Adressendatei:\fbox{\Phone}

\VerbX{Geschlecht}  Ist das Geschlecht in der Adressendatei: \fbox{\Geschlecht}

\VerbX{FreeII}  Ist noch nicht zugewiesen  in der Adressendatei:\fbox{\FreeII}

\VerbX{Comment}  Ist der Kommentar in der Adressendatei: \fbox{\Comment}

\end{labeling}
\closing{Hier kommt die Gru�formel}


}% ende von \Serial

\end{document}
%</exa>
%    \end{macrocode}
%
% \subsection{Example KOMA-Script data file}
%
%    \begin{macrocode}
%<*adr>
\adrentry{Pedersen}{Bjoern}%
{Some Street 1}{11111 Somewhere}%
{11~00~00}{m}{}{Test Kommentar}
\adrentry{Mustermann}{Ingrid}%
{Anotherstreet 99}{99999 Elsewhere}%
{99~99~99}{w}{Hier steht jetzt auch etwas}{Test}
%</adr>
%    \end{macrocode}
% \subsection{Example textmerge data file}
%    \begin{macrocode}
%<*adt>
Pedersen
Bjoern
Some Street 1
11111 Somewhere
11~00~00
m

Test Kommentar
Mustermann
Ingrid
Anotherstreet 99
99999 Elsewhere
99~99~99
w
Hier steht jetz auch etwas
Test
%</adt>
%    \end{macrocode}
%    \begin{macrocode}
%<*del>
Pedersen|Bjoern|Some Street 1|11111 Somewhere|11~00~00|m| |Test Kommentar
Mustermann|Ingrid|Anotherstreet 99|99999 Elsewhere|99~99~99|w|Hier steht jetz auch etwas|Test
%</del>
%    \end{macrocode}
% \subsection{Example configuration file} 
%    \begin{macrocode}
%<*cfg>
%% Some packages which may prove useful: remove  comments as appropriate
%%\RequirePackage[germanb]{babel}
%%\RequirePackage[T1]{fontenc}
%%\RequirePackage{inputenc}
%%\RequirePackage[final]{graphicx}
%%\RequirePackage{times}

%% change certain sizes is allowed here: uncomment and adjust as necessary.
%% The margin will be calculated from the top rim of the paper!
%%\addtolength{\topmargin}{10truemm}
%%\setlength{\headheight}{27truemm}
%%\setlength{\headsep}{4truemm}
%%\setlength{\foldhskip}{6truemm} % adjust for printer
%%\setlength{\footheight}{25truemm}


\signature{Your Name}

\place{Your Place}
\name{Your Name}

\address{%
Your adress name\\%
Your Street\\[1.5ex]%
Your location%
}
\location{\sffamily\scriptsize%
Your Text for the location field
}

\telephone{Your Phone}

\backaddress{\fontsize{7}{8}\selectfont\parbox{70truemm}%
{Your Backaddress,\\ can even be two lines}\rule[-1ex]{0pt}{0pt}}

\DeclareRobustCommand{\MyFirstHead}{%
                \parbox[b]{0.99\textwidth}%
                {%
                \begin{minipage}[b]{98truemm}
                        {\sffamily\bfseries\fontsize{4.1truemm}{6.0truemm}\selectfont
                        \ignorespaces
                         Something for the letter head.}%
                \end{minipage}\hfil
                \begin{minipage}[b]{48truemm}%
                        \hfill\vfill%\parbox[b]{47.5truemm}%
                        {%
                        % maybe ypu want a logo?
                                %\includegraphics[width=47truemm]{logo.eps}%
                        }
                \end{minipage}
                }%
}
\DeclareRobustCommand{\MyFirstFoot}{%
\parbox[t]{\textwidth}
                {\rule{\textwidth}{.4pt}\\[0.6ex]%
                \sffamily\footnotesize%
                \begin{tabular}{lll}
Vorstand:& Sparkasse & Raiffeisenbank \\%
someone &BLZ 700\,530\,70  &BLZ 701\,693\,80\\%
&Kto.-Nr.: 1\,11\,111  &Kto.-Nr.: 999\,999\,999%
                \end{tabular}% 
}% end of parbox
}%

\endinput
%</cfg>
%    \end{macrocode}
%\Finale
